\section{Methodology}\label{sec:method}

This study adopts Systematic Literature Review (SLR) methodology for 
identifying relevant studies and research gaps. This SLR follows 
guidance provided by B. Kitchenham et al. \citep{segress} and
S. Keele et al. \citep{slrguidese}, and incorporates three major
steps: \textbf{Planning the Search} (\ref{subsec:search}), 
\textbf{Study Selection} (\ref{subsec:stusel}),
and \textbf{Data Analysis}.

\subsection{Search Strategy}\label{subsec:search}

As shown in Figure~\ref{fig:stusel},
we employed \textbf{Precise Search Strategy} \citep{optimalsearch}
for article search, i.e., 
we crafted the search string to maximize the relevance 
of resulting studies. More specifically, we followed the following three steps
in order to establish a set of relevant studies:


\begin{enumerate}
 \item Conduct an automated search of studies published within 5 years
  (i.e. after 2022)  based on our crafted search string;
 \item Screen the title and abstract of all articles and filter by
       inclusion/exclusion criteria;
 \item Conduct snowballing search on the result of previous step.
\end{enumerate}

\begin{figure}[ht]

\begin{tikzpicture}[
  every node/.style={fill=white},
  str/.style={rectangle, rounded corners, draw=blue!60, fill=blue!5, very thick, minimum size=10mm},
  db/.style={rectangle, rounded corners, draw=green!60, fill=green!5, very thick, minimum size=10mm},
  filt/.style={rectangle, rounded corners, draw=red!60, fill=red!5, very thick, minimum size=10mm},
  res/.style={rectangle, rounded corners, draw=yellow!60, fill=yellow!5, very thick, minimum size=10mm},
]

  \node[str] (sstr) {
    \begin{tabular} {l}
      \textbf{Search String} \\
       \\
      LLM, \\
      Software Supply Chain, \\
      \\
    \end{tabular}
  };

  \node[db] (sdb) [right=of sstr] {
    \begin{tabular} {ll}
      \multicolumn{2}{l}{\textbf{Search Database}}\\
      ACM & 238 \\
      arXiv & 7 \\
      IEEE & 106 \\
      ScienceDirect & 69 \\
      & \\
    \end{tabular}
  };


  \node[filt] (sfilt) [right=of sdb] {
    \begin{tabular} {ll}
      \multicolumn{2}{l}{\textbf{Apply Selection Criteria}}\\
      \multicolumn{2}{l}{\small (from Table~\ref{tab:criteria})} \\
      ACM & 38 \\
      arXiv & 5 \\
      IEEE & 28 \\
      ScienceDirect & 6 \\
    \end{tabular}
  };

  \node[res] (result) [below=of sdb] {
    \begin{tabular}{l}
      \\
      \textbf{Results in:} \\
      81 Studies \\
      \\
    \end{tabular}
  };

  \node[res] (defense) [below=of result, xshift=-2cm] {
    \begin{tabular}{l}
      \textbf{LLM for Defense} \\
      {\small (Section~\ref{sec:rq2})} \\
      58 Studies \\
    \end{tabular}
  };

  \node[res] (attack) [right=of defense] {
    \begin{tabular}{l}
      \textbf{LLM for Attack} \\
      {\small (Section~\ref{sec:rq3})} \\
      23 Studies \\
    \end{tabular}
  };

  \draw [thick, ->] (sstr.east) -- (sdb.west);
  \draw [thick, ->] (sdb.east) -- (sfilt.west);
  \draw [thick, ->] (sfilt.south) |- node[xshift=-0.5cm] {
    \begin{tabular}{l}
    Snowballing \\
    Search (+3 studies) \\
    \end{tabular}
  } (result.east);
  \draw[thick, ->] (result.south) -- (defense.north);
  \draw[thick, ->] (result.south) -- (attack.north);
\end{tikzpicture}

\caption{The process of study identification and selection.}
\label{fig:stusel}
\end{figure}

\par Our search string needs to combine two sets of keywords, one related
to Software Supply Chain, the other concerned about Large Language Models.
For automated search, intra-set keywords are OR-ed, and inter-set keywords
are AND-ed to form the final search string.
The complete set of keywords are as follows:

\begin{itemize}
  \item \textit{Phrases related to Software Supply Chain}: supply chain, software.

  \item \textit{Phrases related to LLMs}: language model, agent.

  \item \textit{Phrases of task definition}: attack, breach, defend, detect, vulnerability.
\end{itemize}

\par It is important to note that the list of keywords related to 
LLMS we setup includes \textit{agent}, that does not seem to be
necessarily related to LLMs. The reason for this is that we observe 
recent prevalence in LLM-based agentic technologies \citep{agents} (the technique
behind many AI assistant, such as OpenClaw and \texttt{opencode}) ,
and that we want to avoid omitting
articles related to our research as much as possible.

\paragraph*{Search Databases.} After determining the search strings, we
employed the SLR Tool\citep{slrtool} and conducted
an automated search across four widely used databases
\footnote{ACM Digital Library, arXiv, IEEE Xplore, Science Direct}, 
which cover most published
articles and preprints of under-review articles. Given that the first
LLM trained on code corpus was published in 2021\citep{codex}, 
we focus our search on articles
published from that year onward. The search results from each database were
merged and deduplicated with SLR Tool\citep{slrtool}.
As a result, we obtained 238 articles from ACM Digital Library, 
7 articles from arXiv, 106 articles from IEEE Xplore, 
and 69 articles from Science Direct.

\subsection{Study Selection}\label{subsec:stusel}

\par Using our search strategy, we initially obtained 420 articles
that potentially relate to our research. Next, we further 
evaluate the relevance of these articles based on our inclusion
and exclusion criteria, as shown in Table~\ref{tab:criteria}.
Unlike previous SLRs \citep{llm4se2}, we do not exclude papers
published in workshop track, because their reference lists
may benefit our snowballing search.
In this step, we conducted a manual selection by examining the 
title, abstract, and keywords of each article, resulting
in a total of 81 papers.

\begin{table}[ht]
  \centering
  \caption{Inclusion and Exclusion Criteria for Study Selection}
  \label{tab:criteria}
\begin{tabular}{ll}
\thickhline
\multicolumn{2}{c}{\textbf{Inclusion Criteria}} \\ \hline
    (1) & The study claims an LLM is used. \\
    (2) & The study involves Software Supply Chain. \\
    (3) & The study is full-text available. \\
     & \\
\thickhline
\multicolumn{2}{c}{\textbf{Exclusion Criteria}} \\ \hline
    (1) & The study is concerned with LLM supply chain. \\
    (2) & Non-English written literature. \\
    (3) & Tool demos and editorials; \\
     & \\
\thickhline
\thickhline
\end{tabular}
\end{table}

\paragraph*{Snowballing Search}
To identify any additional 
 possibly relevant primary studies, we conducted a snowballing 
 search. Snowballing search refers to using the reference list
 of an article or the citations to the article to identify additional
 articles (which must also comply with our
 selection criteria shown in Table~\ref{tab:criteria}). 
As a result, we obtained an additional 3 studies.


